\begin{explanationbox}
    \textbf{\textcolor{CExplanation}{一句话直觉}}

    \textbf{两条切线虽然切在不同的点，却经过同一个外点 $P$。}
    因此两个切点一定满足同一个一次方程，这个方程就是切点弦方程。

    \medskip
    \textbf{\textcolor{CExplanation}{缩放看见对应关系}}

\begin{center}
\begin{tikzpicture}[x=1cm,y=1cm,font=\small,
  every node/.style={inner sep=3pt},line cap=round]

\begin{scope}[shift={(-2.7,0)},x=1.6cm,y=1.6cm]
\draw[CStatement,very thick] (0,0) circle (1);
\coordinate (N) at (.48,.36);
\draw[orange!85!black,thick] (0,1)--(1.33333,1)--(.96,-.28);
\draw[CTrap,very thick] (0,1)--(.96,-.28);
\draw[dashed,CProof] (0,0)--(1.33333,1);
\draw[CProof] (.54,.405)--(.495,.465)--(.435,.42);
\fill (0,0) circle (1pt) node[below left] {$O$};
\fill (N) circle (1pt) node[below left] {$N$};
\fill (1.33333,1) circle (1pt) node[above] {$Q$};
\node at (0,-1.45) {单位圆：$ON\perp A'B'$};
\node at (0,1.45) {$u^2+v^2=1$};
\node[above left] at (0,1) {$A'$}; \node[below right] at (.96,-.28) {$B'$};
\end{scope}
\node at (.4,-3.2) {$\xrightarrow{\ x=3u,\ y=2v\ }$};
\begin{scope}[shift={(3.5,0)},x=.95cm,y=.95cm]
\draw[CStatement,very thick] (0,0) ellipse [x radius=3,y radius=2];
\draw[orange!85!black,thick] (0,2)--(4,2)--(2.88,-.56);
\draw[CTrap,very thick] (0,2)--(2.88,-.56);
\draw[dashed,CProof] (0,0)--(4,2);
\fill (0,0) circle (1.4pt) node[below left] {$O$};
\fill (1.44,.72) circle (1.4pt) node[below left] {$M$};
\fill (4,2) circle (1.4pt) node[above] {$P$};
\node[above left] at (0,2) {$A$}; \node[below right] at (2.88,-.56) {$B$};
\node at (0,-2.45) {椭圆：$O,M,P$ 共线};
\node at (0,3.2) {$\dfrac{x^2}{9}+\dfrac{y^2}{4}=1$};
\end{scope}

\end{tikzpicture}

{\small 缩放保持切点、直线与中点；圆中的直角一般不会保留。}
\end{center}


    图中取 $a=3,b=2$、$P(4,2)$，对应 $Q(\frac43,1)$。
    把椭圆上每一点按 $u=x/a,v=y/b$ 缩放，就得到单位圆。
    直线与曲线的交点一一对应，因此切线仍对应切线。
    圆中切点的坐标乘积关系 $u_0u_1+v_0v_1=1$（见方法一），
    换回坐标就是 $x_0x_1/a^2+y_0y_1/b^2=1$。

    \medskip
    \textbf{\textcolor{CExplanation}{为什么不必算出两个切点？}}

    已知椭圆在点 $T(x_1,y_1)$ 处的切线方程是
    $\frac{x_1x}{a^2}+\frac{y_1y}{b^2}=1$。
    因为该切线经过 $P$，可代入 $(x_0,y_0)$；交换两个坐标乘积的书写次序后，就发现 $T$ 位于
    $\frac{x_0x}{a^2}+\frac{y_0y}{b^2}=1$ 上。另一个切点同理。
    \textbf{这里只需确定两个切点共同满足的直线方程，而不必先求它们的坐标。}

    \medskip
    \textbf{\textcolor{CExplanation}{一眼区别：切线与切点弦}}
    \begin{itemize}[leftmargin=*, itemsep=0.35em]
        \item 切线 $PA,PB$ 都经过外点 $P$，分别只与椭圆相切于 $A,B$。
        \item 切点弦 $AB$ 经过 $A,B$，\textbf{不经过 $P$}：把 $P$ 代入切点弦的左端，得到 $S>1$，而不是 $1$。
        \item 圆的类似公式为 $x_0x+y_0y=R^2$；椭圆可理解为把圆分别沿两个坐标轴伸缩得到。
    \end{itemize}

    \medskip
    \textbf{\textcolor{CExplanation}{中点为什么沿 $OP$？}}

    缩放后的单位圆中，切点弦是 $u_0u+v_0v=1$，其垂足就是圆弦中点
    $\left(\frac{u_0}{S},\frac{v_0}{S}\right)$，这里 $S=u_0^2+v_0^2$。
    把坐标乘回 $a,b$，就得到椭圆弦中点
    $M\left(\frac{x_0}{S},\frac{y_0}{S}\right)$。
    注意：\textbf{缩放保持线段中点，但一般不保持垂直关系。}

    \medskip
    \textbf{\textcolor{CExplanation}{考试使用信号}}
    \begin{itemize}[leftmargin=*, itemsep=0.3em]
        \item 题目同时出现``椭圆外一点''、``两条切线''、``连接两切点''：优先使用切点弦公式。
        \item 已知一条切点弦 $Ax+By=C$（$C\ne0$）：先化成右端为 $1$ 的形式，再比较系数求外点。
        \item 若题目还涉及弦中点或 $O,P,M$ 共线，可使用 $M=P/S$ 这一补充结论。
    \end{itemize}
\end{explanationbox}
